\documentclass[10pt,a4paper]{amsart}
\usepackage{amsmath,amssymb,amsthm,graphicx}
\usepackage{verbatim}
\newenvironment{reference}{\comment}{\endcomment}
%\newenvironment{reference}{}{}
\newenvironment{slogan}{\comment}{\endcomment}
\newenvironment{history}{\comment}{\endcomment}

% For commutative diagrams we use Xy-pic
\usepackage[all]{xy}

% We use 2cell for 2-commutative diagrams.
\xyoption{2cell}
% Theorem environments.
%
\theoremstyle{plain}
\newtheorem{theorem}[subsection]{Theorem}
\newtheorem{proposition}[subsection]{Proposition}
\newtheorem{lemma}[subsection]{Lemma}

\theoremstyle{definition}
\newtheorem{definition}[subsection]{Definition}
\newtheorem{example}[subsection]{Example}
\newtheorem{exercise}[subsection]{Exercise}
\newtheorem{situation}[subsection]{Situation}

\theoremstyle{remark}
\newtheorem{remark}[subsection]{Remark}
\newtheorem{remarks}[subsection]{Remarks}

\numberwithin{equation}{subsection}

% Macros
%
\def\lim{\mathop{\mathrm{lim}}\nolimits}
\def\colim{\mathop{\mathrm{colim}}\nolimits}
\def\Spec{\mathop{\mathrm{Spec}}}
\def\Hom{\mathop{\mathrm{Hom}}\nolimits}
\def\Ext{\mathop{\mathrm{Ext}}\nolimits}
\def\SheafHom{\mathop{\mathcal{H}\!\mathit{om}}\nolimits}
\def\SheafExt{\mathop{\mathcal{E}\!\mathit{xt}}\nolimits}
\def\Sch{\mathit{Sch}}
\def\Mor{\mathop{\mathrm{Mor}}\nolimits}
\def\Ob{\mathop{\mathrm{Ob}}\nolimits}
\def\Sh{\mathop{\mathit{Sh}}\nolimits}
\def\NL{\mathop{N\!L}\nolimits}
\def\CH{\mathop{\mathrm{CH}}\nolimits}
\def\proetale{{pro\text{-}\acute{e}tale}}
\def\etale{{\acute{e}tale}}
\def\QCoh{\mathit{QCoh}}
\def\Ker{\mathop{\mathrm{Ker}}}
\def\Im{\mathop{\mathrm{Im}}}
\def\Coker{\mathop{\mathrm{Coker}}}
\def\Coim{\mathop{\mathrm{Coim}}}

% Boxtimes
%
\DeclareMathSymbol{\boxtimes}{\mathbin}{AMSa}{"02}

%
% Macros for moduli stacks/spaces
%
\def\QCohstack{\mathcal{QC}\!\mathit{oh}}
\def\Cohstack{\mathcal{C}\!\mathit{oh}}
\def\Spacesstack{\mathcal{S}\!\mathit{paces}}
\def\Quotfunctor{\mathrm{Quot}}
\def\Hilbfunctor{\mathrm{Hilb}}
\def\Curvesstack{\mathcal{C}\!\mathit{urves}}
\def\Polarizedstack{\mathcal{P}\!\mathit{olarized}}
\def\Complexesstack{\mathcal{C}\!\mathit{omplexes}}
% \Pic is the operator that assigns to X its picard group, usage \Pic(X)
% \Picardstack_{X/B} denotes the Picard stack of X over B
% \Picardfunctor_{X/B} denotes the Picard functor of X over B
\def\Pic{\mathop{\mathrm{Pic}}\nolimits}
\def\Picardstack{\mathcal{P}\!\mathit{ic}}
\def\Picardfunctor{\mathrm{Pic}}
\def\Deformationcategory{\mathcal{D}\!\mathit{ef}}
\usepackage{cleveref}
\crefname{originallemma}{original lemma}{original lemmas}
\Crefname{originallemma}{Original lemma}{Original lemmas}
\makeatletter
\newlabel{fields-lemma-primitive-element}{{030N}{source}{}{sourcepointer-1}{}}
\newlabel{fields-lemma-primitive-element@cref}{{[originallemma][0][]030N}{[1][0][]source}{}{}{}}
\makeatother
\crefname{originalsection}{original section}{original sections}
\Crefname{originalsection}{Original section}{Original sections}
\makeatletter
\newlabel{fields-section-separable-extensions}{{09GZ}{source}{}{sourcepointer-2}{}}
\newlabel{fields-section-separable-extensions@cref}{{[originalsection][0][]09GZ}{[1][0][]source}{}{}{}}
\makeatother
\title{For a separable algebraic extension the multiplication map is a localisation of the square}
\author{The Stacks Project Authors}
\date{stacks-project:algebra.tex:local-snapshot-20261012; GNU FDL 1.2 or later}
\begin{document}
\maketitle
\noindent\textit{Editorial note.} All delivered mathematical text is the author's own, copied from the pinned Stacks Project chapter file algebra.tex; no statement or proof step was rewritten, reordered or removed.\par\medskip
\noindent\textit{Editorial note.} Preamble provenance, as required for this source: the frozen snapshot stages the chapter's own preamble file beside it (preamble.tex, 217 lines, 8116 bytes, the same shared Stacks-wide preamble the published chow and etale-cohomology cards use); the first line of the pinned chapter file is \textbackslash{}input\{preamble\}, so no preamble text lives inside the chapter itself. The card takes that file's lines 8-18 (the verbatim package and the three \textbackslash{}comment-based environments reference/slogan/history at lines 9-12, section 13.2, together with the Xy-pic setup) and its lines 154 to the runtime end of the file (the native theorem declarations and mathematical macros). Not one character was changed, no package was added, no macro was defined, and the excerpt serves only as preamble, never as a span coordinate.\par\medskip
\noindent\textit{Editorial note.} The same-chapter pointers are located at their \textbackslash{}label lines inside the pinned chapter file; the cross-chapter pointers are located in the project's own tag list staged in the snapshot (sources/stacks-chow/tags.txt), at the verbatim lines carrying their canonical tags.\par\medskip
\noindent\textit{Editorial note.} The counted statement and proof refer to results of the Fields chapter by label, one of them a section; they are recorded as navigation stubs with the appropriate kinds, and the arguments using them stay inside the counted span.\par\medskip
\noindent\textit{Editorial note.} The counted spans contain no \textbackslash{}cite, so the card carries no bibliography.\par\medskip
\noindent\textit{Allowed standard prerequisites (editorial).}\begin{itemize}
\item The primitive element theorem for finite separable extensions. Invoked in the counted proof at line 10447. Writes a finite separable extension as generated by one element with a separable minimal polynomial.
\item Colimits and exactness of localisation. Used in the second half of the counted proof (lines 10460-10471). Passes the finite case to arbitrary separable algebraic extensions and shows the kernel is killed.
\end{itemize}\medskip
\ifdefined\hypertarget\hypertarget{sourcepointer-1}{}\fi
\noindent\textit{Source reference.} Original lemma 030N: sources/stacks-chow/tags.txt, source lines 3954--3954. Navigation pointer for \textbackslash{}ref\{fields-lemma-primitive-element\} used inside the counted span (line 10447); the argument that consumes it is retained in the span.\par\medskip
\ifdefined\hypertarget\hypertarget{sourcepointer-2}{}\fi
\noindent\textit{Source reference.} Original section 09GZ: sources/stacks-chow/tags.txt, source lines 11987--11987. Navigation pointer for \textbackslash{}ref\{fields-section-separable-extensions\} used inside the counted span (line 10450); the argument that consumes it is retained in the span.\par\medskip
\begin{lemma}
\label{lemma-separable-algebraic-diagonal}
Let $k'/k$ be a separable algebraic extension.
Then there exists a multiplicative subset $S \subset k' \otimes_k k'$
such that the multiplication map $k' \otimes_k k' \to k'$
is identified with $k' \otimes_k k' \to S^{-1}(k' \otimes_k k')$.
\end{lemma}

\begin{proof}
First assume $k'/k$ is finite separable. Then $k' = k(\alpha)$,
see Fields, Lemma \ref{fields-lemma-primitive-element}.
Let $P \in k[x]$ be the minimal polynomial of $\alpha$ over $k$.
Then $P$ is an irreducible, separable, monic polynomial, see
Fields, Section \ref{fields-section-separable-extensions}.
Then $k'[x]/(P) \to k' \otimes_k k'$,
$\sum \alpha_i x^i \mapsto \alpha_i \otimes \alpha^i$ is an isomorphism.
We can factor $P = (x - \alpha) Q$ in $k'[x]$ and since $P$
is separable we see that $Q(\alpha) \not = 0$.
Then it is clear that the multiplicative set $S'$ generated by
$Q$ in $k'[x]/(P)$ works, i.e., that $k' = (S')^{-1}(k'[x]/(P))$.
By transport of structure the image $S$ of $S'$ in $k' \otimes_k k'$
works.

\medskip\noindent
In the general case we write $k' = \bigcup k_i$ as the union
of its finite subfield extensions over $k$. For each $i$ there
is a multiplicative subset $S_i \subset k_i \otimes_k k_i$
such that $k_i = S_i^{-1}(k_i \otimes_k k_i)$. Let
$S$ be the multiplicative closure of $\bigcup S_i \subset k' \otimes_k k'$.
Clearly, the multiplication maps sends every element of $S$
to an invertible element of $k'$. Using that $k' \otimes_k k'$
is the union of the rings $k_i \otimes_k k_i$ we see that
every element in the kernel of the multiplication map is mapped to zero
in $S^{-1}(k' \otimes_k k')$. Using exactness of localization
the result follows.
\end{proof}

\end{document}
